\documentclass[10pt,a4paper]{amsart}
\usepackage{amsmath,amssymb,amsthm,graphicx}
% For commutative diagrams we use Xy-pic
\usepackage[all]{xy}

% We use 2cell for 2-commutative diagrams.
\xyoption{2cell}
% Theorem environments.
%
\theoremstyle{plain}
\newtheorem{theorem}[subsection]{Theorem}
\newtheorem{proposition}[subsection]{Proposition}
\newtheorem{lemma}[subsection]{Lemma}

\theoremstyle{definition}
\newtheorem{definition}[subsection]{Definition}
\newtheorem{example}[subsection]{Example}
\newtheorem{exercise}[subsection]{Exercise}
\newtheorem{situation}[subsection]{Situation}

\theoremstyle{remark}
\newtheorem{remark}[subsection]{Remark}
\newtheorem{remarks}[subsection]{Remarks}

\numberwithin{equation}{subsection}

% Macros
%
\def\lim{\mathop{\mathrm{lim}}\nolimits}
\def\colim{\mathop{\mathrm{colim}}\nolimits}
\def\Spec{\mathop{\mathrm{Spec}}}
\def\Hom{\mathop{\mathrm{Hom}}\nolimits}
\def\Ext{\mathop{\mathrm{Ext}}\nolimits}
\def\SheafHom{\mathop{\mathcal{H}\!\mathit{om}}\nolimits}
\def\SheafExt{\mathop{\mathcal{E}\!\mathit{xt}}\nolimits}
\def\Sch{\mathit{Sch}}
\def\Mor{\mathop{\mathrm{Mor}}\nolimits}
\def\Ob{\mathop{\mathrm{Ob}}\nolimits}
\def\Sh{\mathop{\mathit{Sh}}\nolimits}
\def\NL{\mathop{N\!L}\nolimits}
\def\CH{\mathop{\mathrm{CH}}\nolimits}
\def\proetale{{pro\text{-}\acute{e}tale}}
\def\etale{{\acute{e}tale}}
\def\QCoh{\mathit{QCoh}}
\def\Ker{\mathop{\mathrm{Ker}}}
\def\Im{\mathop{\mathrm{Im}}}
\def\Coker{\mathop{\mathrm{Coker}}}
\def\Coim{\mathop{\mathrm{Coim}}}

% Boxtimes
%
\DeclareMathSymbol{\boxtimes}{\mathbin}{AMSa}{"02}

%
% Macros for moduli stacks/spaces
%
\def\QCohstack{\mathcal{QC}\!\mathit{oh}}
\def\Cohstack{\mathcal{C}\!\mathit{oh}}
\def\Spacesstack{\mathcal{S}\!\mathit{paces}}
\def\Quotfunctor{\mathrm{Quot}}
\def\Hilbfunctor{\mathrm{Hilb}}
\def\Curvesstack{\mathcal{C}\!\mathit{urves}}
\def\Polarizedstack{\mathcal{P}\!\mathit{olarized}}
\def\Complexesstack{\mathcal{C}\!\mathit{omplexes}}
% \Pic is the operator that assigns to X its picard group, usage \Pic(X)
% \Picardstack_{X/B} denotes the Picard stack of X over B
% \Picardfunctor_{X/B} denotes the Picard functor of X over B
\def\Pic{\mathop{\mathrm{Pic}}\nolimits}
\def\Picardstack{\mathcal{P}\!\mathit{ic}}
\def\Picardfunctor{\mathrm{Pic}}
\def\Deformationcategory{\mathcal{D}\!\mathit{ef}}
\newenvironment{reference}{}{}
\newenvironment{slogan}{}{}
\newenvironment{history}{}{}
\usepackage{cleveref}
\crefname{originallemma}{original lemma}{original lemmas}
\Crefname{originallemma}{Original lemma}{Original lemmas}
\makeatletter
\newlabel{more-algebra-lemma-K-flat-two-out-of-three}{{06Y2}{source}{}{sourcepointer-1}{}}
\newlabel{more-algebra-lemma-K-flat-two-out-of-three@cref}{{[originallemma][0][]06Y2}{[1][0][]source}{}{}{}}
\makeatother
\crefname{originallemma}{original lemma}{original lemmas}
\Crefname{originallemma}{Original lemma}{Original lemmas}
\makeatletter
\newlabel{cohomology-lemma-alternating-usual}{{01FM}{source}{}{sourcepointer-2}{}}
\newlabel{cohomology-lemma-alternating-usual@cref}{{[originallemma][0][]01FM}{[1][0][]source}{}{}{}}
\makeatother
\crefname{originallemma}{original lemma}{original lemmas}
\Crefname{originallemma}{Original lemma}{Original lemmas}
\makeatletter
\newlabel{more-algebra-lemma-derived-tor-quasi-isomorphism}{{064K}{source}{}{sourcepointer-3}{}}
\newlabel{more-algebra-lemma-derived-tor-quasi-isomorphism@cref}{{[originallemma][0][]064K}{[1][0][]source}{}{}{}}
\makeatother
\crefname{originallemma}{original lemma}{original lemmas}
\Crefname{originallemma}{Original lemma}{Original lemmas}
\makeatletter
\newlabel{varieties-lemma-kunneth}{{0BED}{source}{}{sourcepointer-4}{}}
\newlabel{varieties-lemma-kunneth@cref}{{[originallemma][0][]0BED}{[1][0][]source}{}{}{}}
\makeatother
\crefname{originalsection}{original section}{original sections}
\Crefname{originalsection}{Original section}{Original sections}
\makeatletter
\newlabel{cohomology-section-cech-cohomology-of-complexes}{{01FP}{source}{}{sourcepointer-5}{}}
\newlabel{cohomology-section-cech-cohomology-of-complexes@cref}{{[originalsection][0][]01FP}{[1][0][]source}{}{}{}}
\makeatother
\crefname{originalequation}{original equation}{original equations}
\Crefname{originalequation}{Original equation}{Original equations}
\makeatletter
\newlabel{cohomology-equation-needs-signs}{{07MB}{source}{}{sourcepointer-6}{}}
\newlabel{cohomology-equation-needs-signs@cref}{{[originalequation][0][]07MB}{[1][0][]source}{}{}{}}
\makeatother
\crefname{originallemma}{original lemma}{original lemmas}
\Crefname{originallemma}{Original lemma}{Original lemmas}
\makeatletter
\newlabel{cohomology-lemma-diagrams-commute}{{0FKV}{source}{}{sourcepointer-7}{}}
\newlabel{cohomology-lemma-diagrams-commute@cref}{{[originallemma][0][]0FKV}{[1][0][]source}{}{}{}}
\makeatother
\crefname{originallemma}{original lemma}{original lemmas}
\Crefname{originallemma}{Original lemma}{Original lemmas}
\makeatletter
\newlabel{coherent-lemma-quasi-coherent-affine-cohomology-zero}{{01XB}{source}{}{sourcepointer-8}{}}
\newlabel{coherent-lemma-quasi-coherent-affine-cohomology-zero@cref}{{[originallemma][0][]01XB}{[1][0][]source}{}{}{}}
\makeatother
\crefname{originallemma}{original lemma}{original lemmas}
\Crefname{originallemma}{Original lemma}{Original lemmas}
\makeatletter
\newlabel{cohomology-lemma-cech-complex-complex-computes}{{0FLH}{source}{}{sourcepointer-9}{}}
\newlabel{cohomology-lemma-cech-complex-complex-computes@cref}{{[originallemma][0][]0FLH}{[1][0][]source}{}{}{}}
\makeatother
\crefname{originallemma}{original lemma}{original lemmas}
\Crefname{originallemma}{Original lemma}{Original lemmas}
\makeatletter
\newlabel{morphisms-lemma-affine-permanence}{{01SG}{source}{}{sourcepointer-10}{}}
\newlabel{morphisms-lemma-affine-permanence@cref}{{[originallemma][0][]01SG}{[1][0][]source}{}{}{}}
\makeatother
\crefname{originallemma}{original lemma}{original lemmas}
\Crefname{originallemma}{Original lemma}{Original lemmas}
\makeatletter
\newlabel{morphisms-lemma-pushforward-flat-affine}{{0FLM}{source}{}{sourcepointer-11}{}}
\newlabel{morphisms-lemma-pushforward-flat-affine@cref}{{[originallemma][0][]0FLM}{[1][0][]source}{}{}{}}
\makeatother
\crefname{originallemma}{original lemma}{original lemmas}
\Crefname{originallemma}{Original lemma}{Original lemmas}
\makeatletter
\newlabel{coherent-lemma-relative-affine-cohomology}{{089W}{source}{}{sourcepointer-12}{}}
\newlabel{coherent-lemma-relative-affine-cohomology@cref}{{[originallemma][0][]089W}{[1][0][]source}{}{}{}}
\makeatother
\crefname{originallemma}{original lemma}{original lemmas}
\Crefname{originallemma}{Original lemma}{Original lemmas}
\makeatletter
\newlabel{coherent-lemma-separated-case-relative-cech}{{01XL}{source}{}{sourcepointer-13}{}}
\newlabel{coherent-lemma-separated-case-relative-cech@cref}{{[originallemma][0][]01XL}{[1][0][]source}{}{}{}}
\makeatother
\crefname{originallemma}{original lemma}{original lemmas}
\Crefname{originallemma}{Original lemma}{Original lemmas}
\makeatletter
\newlabel{lemma-cohomology-base-change}{{08EU}{source}{}{sourcepointer-14}{}}
\newlabel{lemma-cohomology-base-change@cref}{{[originallemma][0][]08EU}{[1][0][]source}{}{}{}}
\makeatother
\crefname{originalsection}{original section}{original sections}
\Crefname{originalsection}{Original section}{Original sections}
\makeatletter
\newlabel{varieties-section-kunneth}{{0BEC}{source}{}{sourcepointer-15}{}}
\newlabel{varieties-section-kunneth@cref}{{[originalsection][0][]0BEC}{[1][0][]source}{}{}{}}
\makeatother
\crefname{originalsection}{original section}{original sections}
\Crefname{originalsection}{Original section}{Original sections}
\makeatletter
\newlabel{cohomology-remark-cup-product}{{0B68}{source}{}{sourcepointer-16}{}}
\newlabel{cohomology-remark-cup-product@cref}{{[originalsection][0][]0B68}{[1][0][]source}{}{}{}}
\makeatother
\crefname{originalsection}{original section}{original sections}
\Crefname{originalsection}{Original section}{Original sections}
\makeatletter
\newlabel{cohomology-section-cup-product}{{0FKU}{source}{}{sourcepointer-17}{}}
\newlabel{cohomology-section-cup-product@cref}{{[originalsection][0][]0FKU}{[1][0][]source}{}{}{}}
\makeatother
\title{Derived Kunneth isomorphism for base-flat quasi-coherent sheaves over an affine base}
\author{The Stacks Project Authors}
\date{stacks-project:perfect.tex:local-snapshot-20261011; GNU FDL 1.2 or later}
\begin{document}
\maketitle
\noindent\textit{Editorial note.} No properness or Noetherian condition is added toX,Y,S. The source assumesX,Y quasi-compact with affine diagonal andF,G flat overS.\par\medskip
\noindent\textit{Editorial note.} The ordinary field case invarieties.tex is not separately counted.\par\medskip
\noindent\textit{Allowed standard prerequisites (editorial).}\begin{itemize}
\item Explicit prerequisite boundary. Quasi-coherent sheaves on schemes with affine diagonal, finite affine covers, cup products andCech resolutions, K-flat tensor products, Leray and flat cohomology base change. The localK-flat argument is retained.
\end{itemize}\medskip
\ifdefined\hypertarget\hypertarget{sourcepointer-1}{}\fi
\noindent\textit{Source reference.} Original lemma 06Y2: focused-seeds/AGP/sources/stacks-snapshot-20261011/tags.txt, source lines 8871--8871. Stable Stacks tag for an expressly allowed foundational or notational input; the counted construction remains the two concrete core obligations. The retained local core is not replaced by this pointer. See standard\_prerequisites.\par\medskip
\ifdefined\hypertarget\hypertarget{sourcepointer-2}{}\fi
\noindent\textit{Source reference.} Original lemma 01FM: focused-seeds/AGP/sources/stacks-snapshot-20261011/tags.txt, source lines 1875--1875. Stable Stacks tag for an expressly allowed foundational or notational input; the counted construction remains the two concrete core obligations. The retained local core is not replaced by this pointer. See standard\_prerequisites.\par\medskip
\ifdefined\hypertarget\hypertarget{sourcepointer-3}{}\fi
\noindent\textit{Source reference.} Original lemma 064K: focused-seeds/AGP/sources/stacks-snapshot-20261011/tags.txt, source lines 7869--7869. Stable Stacks tag for an expressly allowed foundational or notational input; the counted construction remains the two concrete core obligations. The retained local core is not replaced by this pointer. See standard\_prerequisites.\par\medskip
\ifdefined\hypertarget\hypertarget{sourcepointer-4}{}\fi
\noindent\textit{Source reference.} Original lemma 0BED: focused-seeds/AGP/sources/stacks-snapshot-20261011/tags.txt, source lines 14346--14346. Stable Stacks tag for an expressly allowed foundational or notational input; the counted construction remains the two concrete core obligations. The retained local core is not replaced by this pointer. See standard\_prerequisites.\par\medskip
\ifdefined\hypertarget\hypertarget{sourcepointer-5}{}\fi
\noindent\textit{Source reference.} Original section 01FP: focused-seeds/AGP/sources/stacks-snapshot-20261011/tags.txt, source lines 1878--1878. Stable Stacks tag for an expressly allowed foundational or notational input; the counted construction remains the two concrete core obligations. The retained local core is not replaced by this pointer. See standard\_prerequisites.\par\medskip
\ifdefined\hypertarget\hypertarget{sourcepointer-6}{}\fi
\noindent\textit{Source reference.} Original equation 07MB: focused-seeds/AGP/sources/stacks-snapshot-20261011/tags.txt, source lines 9721--9721. Stable Stacks tag for an expressly allowed foundational or notational input; the counted construction remains the two concrete core obligations. The retained local core is not replaced by this pointer. See standard\_prerequisites.\par\medskip
\ifdefined\hypertarget\hypertarget{sourcepointer-7}{}\fi
\noindent\textit{Source reference.} Original lemma 0FKV: focused-seeds/AGP/sources/stacks-snapshot-20261011/tags.txt, source lines 19475--19475. Stable Stacks tag for an expressly allowed foundational or notational input; the counted construction remains the two concrete core obligations. The retained local core is not replaced by this pointer. See standard\_prerequisites.\par\medskip
\ifdefined\hypertarget\hypertarget{sourcepointer-8}{}\fi
\noindent\textit{Source reference.} Original lemma 01XB: focused-seeds/AGP/sources/stacks-snapshot-20261011/tags.txt, source lines 2512--2512. Stable Stacks tag for an expressly allowed foundational or notational input; the counted construction remains the two concrete core obligations. The retained local core is not replaced by this pointer. See standard\_prerequisites.\par\medskip
\ifdefined\hypertarget\hypertarget{sourcepointer-9}{}\fi
\noindent\textit{Source reference.} Original lemma 0FLH: focused-seeds/AGP/sources/stacks-snapshot-20261011/tags.txt, source lines 19497--19497. Stable Stacks tag for an expressly allowed foundational or notational input; the counted construction remains the two concrete core obligations. The retained local core is not replaced by this pointer. See standard\_prerequisites.\par\medskip
\ifdefined\hypertarget\hypertarget{sourcepointer-10}{}\fi
\noindent\textit{Source reference.} Original lemma 01SG: focused-seeds/AGP/sources/stacks-snapshot-20261011/tags.txt, source lines 2337--2337. Stable Stacks tag for an expressly allowed foundational or notational input; the counted construction remains the two concrete core obligations. The retained local core is not replaced by this pointer. See standard\_prerequisites.\par\medskip
\ifdefined\hypertarget\hypertarget{sourcepointer-11}{}\fi
\noindent\textit{Source reference.} Original lemma 0FLM: focused-seeds/AGP/sources/stacks-snapshot-20261011/tags.txt, source lines 19502--19502. Stable Stacks tag for an expressly allowed foundational or notational input; the counted construction remains the two concrete core obligations. The retained local core is not replaced by this pointer. See standard\_prerequisites.\par\medskip
\ifdefined\hypertarget\hypertarget{sourcepointer-12}{}\fi
\noindent\textit{Source reference.} Original lemma 089W: focused-seeds/AGP/sources/stacks-snapshot-20261011/tags.txt, source lines 10511--10511. Stable Stacks tag for an expressly allowed foundational or notational input; the counted construction remains the two concrete core obligations. The retained local core is not replaced by this pointer. See standard\_prerequisites.\par\medskip
\ifdefined\hypertarget\hypertarget{sourcepointer-13}{}\fi
\noindent\textit{Source reference.} Original lemma 01XL: focused-seeds/AGP/sources/stacks-snapshot-20261011/tags.txt, source lines 2522--2522. Stable Stacks tag for an expressly allowed foundational or notational input; the counted construction remains the two concrete core obligations. The retained local core is not replaced by this pointer. See standard\_prerequisites.\par\medskip
\ifdefined\hypertarget\hypertarget{sourcepointer-14}{}\fi
\noindent\textit{Source reference.} Original lemma 08EU: focused-seeds/AGP/sources/stacks-snapshot-20261011/tags.txt, source lines 10685--10685. Stable Stacks tag for an expressly allowed foundational or notational input; the counted construction remains the two concrete core obligations. The retained local core is not replaced by this pointer. See standard\_prerequisites.\par\medskip
\ifdefined\hypertarget\hypertarget{sourcepointer-15}{}\fi
\noindent\textit{Source reference.} Original section 0BEC: focused-seeds/AGP/sources/stacks-snapshot-20261011/tags.txt, source lines 14345--14345. Stable Stacks tag for an expressly allowed foundational or notational input; the counted construction remains the two concrete core obligations. The retained local core is not replaced by this pointer. See standard\_prerequisites.\par\medskip
\ifdefined\hypertarget\hypertarget{sourcepointer-16}{}\fi
\noindent\textit{Source reference.} Original section 0B68: focused-seeds/AGP/sources/stacks-snapshot-20261011/tags.txt, source lines 14061--14061. Stable Stacks tag for an expressly allowed foundational or notational input; the counted construction remains the two concrete core obligations. The retained local core is not replaced by this pointer. See standard\_prerequisites.\par\medskip
\ifdefined\hypertarget\hypertarget{sourcepointer-17}{}\fi
\noindent\textit{Source reference.} Original section 0FKU: focused-seeds/AGP/sources/stacks-snapshot-20261011/tags.txt, source lines 19474--19474. Stable Stacks tag for an expressly allowed foundational or notational input; the counted construction remains the two concrete core obligations. The retained local core is not replaced by this pointer. See standard\_prerequisites.\par\medskip
\begin{lemma}
\label{lemma-kunneth-single-sheaf}
In the situation above the map (\ref{equation-kunneth-single-sheaves}) is an
isomorphism if $S$ is affine, $\mathcal{F}$ and $\mathcal{G}$ are $S$-flat and
quasi-coherent and $X$ and $Y$ are quasi-compact with affine diagonal.
\end{lemma}

\section{K\"unneth formula, II}
\label{section-kunneth}

\noindent
For the case where the base is a field, please see
Varieties, Section \ref{varieties-section-kunneth}.
Consider a cartesian diagram of schemes
$$
\xymatrix{
& X \times_S Y \ar[ld]^p \ar[rd]_q \ar[dd]^f \\
X \ar[rd]_a & & Y \ar[ld]^b \\
& S
}
$$
Let $K \in D(\mathcal{O}_X)$ and $M \in D(\mathcal{O}_Y)$.
There is a canonical map
\begin{equation}
\label{equation-kunneth}
Ra_*K \otimes_{\mathcal{O}_S}^\mathbf{L} Rb_*M
\longrightarrow
Rf_*(Lp^*K \otimes_{\mathcal{O}_{X \times_S Y}}^\mathbf{L} Lq^*M)
\end{equation}
Namely, we can use the maps
$Ra_*K \to Ra_*Rp_* Lp^*K = Rf_*Lp^*K$ and
$Rb_*M \to Rb_*Rq_* Lq^*M = Rf_*Lq^*M$ and then we can use the
relative cup product (Cohomology, Remark \ref{cohomology-remark-cup-product}).

\medskip\noindent
Set $A = \Gamma(S, \mathcal{O}_S)$. There is a global K\"unneth map
\begin{equation}
\label{equation-kunneth-global}
R\Gamma(X, K) \otimes_A^\mathbf{L} R\Gamma(Y, M)
\longrightarrow
R\Gamma(X \times_S Y,
Lp^*K \otimes_{\mathcal{O}_{X \times_S Y}}^\mathbf{L} Lq^*M)
\end{equation}
in $D(A)$. This map is constructed using the pullback maps
$R\Gamma(X, K) \to R\Gamma(X \times_S Y, Lp^*K)$ and
$R\Gamma(Y, M) \to R\Gamma(X \times_S Y, Lq^*M)$ and
the cup product constructed in
Cohomology, Section \ref{cohomology-section-cup-product}.

\noindent
Let $X, Y, S, a, b, p, q, f$ be as in the introduction to this section.
Let $\mathcal{F}$ be an $\mathcal{O}_X$-module.
Let $\mathcal{G}$ be an $\mathcal{O}_Y$-module.
Set $A = \Gamma(S, \mathcal{O}_S)$.
Consider the map
\begin{equation}
\label{equation-kunneth-single-sheaves}
R\Gamma(X, \mathcal{F})
\otimes_A^\mathbf{L}
R\Gamma(Y, \mathcal{G})
\longrightarrow
R\Gamma(X \times_S Y,
p^*\mathcal{F}
\otimes_{\mathcal{O}_{X \times_S Y}}
q^*\mathcal{G})
\end{equation}
in $D(A)$. This map is constructed using the pullback maps
$R\Gamma(X, \mathcal{F}) \to R\Gamma(X \times_S Y, p^*\mathcal{F})$ and
$R\Gamma(Y, \mathcal{G}) \to R\Gamma(X \times_S Y, q^*\mathcal{G})$,
the cup product constructed in
Cohomology, Section \ref{cohomology-section-cup-product}, and
the canonical map
$p^*\mathcal{F}
\otimes_{\mathcal{O}_{X \times_S Y}}^\mathbf{L}
q^*\mathcal{G}
\to
p^*\mathcal{F}
\otimes_{\mathcal{O}_{X \times_S Y}}
q^*\mathcal{G}$.

\begin{lemma}
\label{lemma-K-flat}
Let $f : X \to Y$ be a morphism of schemes with $Y = \Spec(A)$ affine.
Let $\mathcal{U} : X = \bigcup_{i \in I} U_i$ be a finite affine open covering
such that all the finite intersections
$U_{i_0 \ldots i_p} = U_{i_0} \cap \ldots \cap U_{i_p}$
are affine. Let $\mathcal{F}^\bullet$ be a bounded complex of
$f^{-1}\mathcal{O}_Y$-modules. Assume for all $n \in \mathbf{Z}$
the sheaf $\mathcal{F}^n$ is a flat $f^{-1}\mathcal{O}_Y$-module and
$\mathcal{F}^n$ has the structure of a quasi-coherent $\mathcal{O}_X$-module
compatible with the given $p^{-1}\mathcal{O}_Y$-module structure (but the
differentials in the complex $\mathcal{F}^\bullet$ need not
be $\mathcal{O}_X$-linear). Then the complex
$\text{Tot}(\check{\mathcal{C}}^\bullet(\mathcal{U}, \mathcal{F}^\bullet))$
is K-flat as a complex of $A$-modules.
\end{lemma}

\begin{proof}
We may write
$$
\mathcal{F}^\bullet = (\mathcal{F}^a \to \ldots \to \mathcal{F}^b)
$$
Arguing by induction on $b - a$ and considering the distinguished triangle
$$
\mathcal{F}^b[-b] \to (\mathcal{F}^a \to \ldots \to \mathcal{F}^b)
\to (\mathcal{F}^a \to \ldots \to \mathcal{F}^{b - 1}) \to
\mathcal{F}^b[-b + 1]
$$
and using
More on Algebra, Lemma \ref{more-algebra-lemma-K-flat-two-out-of-three}
we reduce to the case where $\mathcal{F}^\bullet$ consists of a
single quasi-coherent $\mathcal{O}_X$-module $\mathcal{F}$
placed in degree $0$. In this case the {\v C}ech complex
for $\mathcal{F}$ and $\mathcal{U}$ is homotopy equivalent to the
alternating {\v C}ech complex, see
Cohomology, Lemma \ref{cohomology-lemma-alternating-usual}.
Since $U_{i_0 \ldots i_p}$ is always affine, we see that
$\mathcal{F}(U_{i_0 \ldots i_p})$ is $A$-flat.
Hence
$\check{\mathcal{C}}_{alt}^\bullet(\mathcal{U}, \mathcal{F})$
is a bounded complex of flat $A$-modules and hence K-flat
by More on Algebra, Lemma
\ref{more-algebra-lemma-derived-tor-quasi-isomorphism}.
\end{proof}

\begin{proof}
We strongly urge the reader to read the proof of
Varieties, Lemma \ref{varieties-lemma-kunneth} first.
Choose finite affine open coverings
$\mathcal{U} : X = \bigcup_{i \in I} U_i$ and
$\mathcal{V} : Y = \bigcup_{j \in J} V_j$.
This determines an affine open covering
$\mathcal{W} : X \times_S Y = \bigcup_{(i, j) \in I \times J} U_i \times_S V_j$.
Note that $\mathcal{W}$ is a refinement of
$\text{pr}_1^{-1}\mathcal{U}$ and of $\text{pr}_2^{-1}\mathcal{V}$.
Thus by the discussion in Cohomology, Section
\ref{cohomology-section-cech-cohomology-of-complexes}
we obtain maps
$$
\check{\mathcal{C}}^\bullet(\mathcal{U}, \mathcal{F})
\to
\check{\mathcal{C}}^\bullet(\mathcal{W}, p^*\mathcal{F})
\quad\text{and}\quad
\check{\mathcal{C}}^\bullet(\mathcal{V}, \mathcal{G})
\to
\check{\mathcal{C}}^\bullet(\mathcal{W}, q^*\mathcal{G})
$$
well defined up to homotopy and compatible with pullback maps on cohomology.
In Cohomology, Equation (\ref{cohomology-equation-needs-signs})
we have constructed a map of complexes
$$
\text{Tot}(
\check{\mathcal{C}}^\bullet(\mathcal{W}, p^*\mathcal{F})
\otimes_A
\check{\mathcal{C}}^\bullet(\mathcal{W}, q^*\mathcal{G}))
\longrightarrow
\check{\mathcal{C}}^\bullet(\mathcal{W},
p^*\mathcal{F} \otimes_{\mathcal{O}_{X \times_S Y}}
q^*\mathcal{G})
$$
which is compatible with the cup product on cohomology by
Cohomology, Lemma \ref{cohomology-lemma-diagrams-commute}.
Combining the above we obtain a map of complexes
\begin{equation}
\label{equation-kunneth-on-cech}
\text{Tot}(
\check{\mathcal{C}}^\bullet(\mathcal{U}, \mathcal{F})
\otimes_A
\check{\mathcal{C}}^\bullet(\mathcal{V}, \mathcal{G}))
\to
\check{\mathcal{C}}^\bullet(\mathcal{W},
p^*\mathcal{F}
\otimes_{\mathcal{O}_{X \times_S Y}}
q^*\mathcal{G})
\end{equation}
We claim this is the map in the statement of the lemma, i.e.,
the source and target of this arrow are the same as the source
and target of (\ref{equation-kunneth-single-sheaves}). Namely, by
Cohomology of Schemes, Lemma
\ref{coherent-lemma-quasi-coherent-affine-cohomology-zero}
and
Cohomology, Lemma \ref{cohomology-lemma-cech-complex-complex-computes}
the canonical maps
$$
\check{\mathcal{C}}^\bullet(\mathcal{U}, \mathcal{F})
\to
R\Gamma(X, \mathcal{F}),
\quad
\check{\mathcal{C}}^\bullet(\mathcal{V}, \mathcal{G})
\to
R\Gamma(Y, \mathcal{G})
$$
and
$$
\check{\mathcal{C}}^\bullet(\mathcal{W},
p^*\mathcal{F} \otimes_{\mathcal{O}_{X \times_S Y}} q^*\mathcal{G})
\to
R\Gamma(X \times_S Y,
p^*\mathcal{F} \otimes_{\mathcal{O}_{X \times_S Y}}
q^*\mathcal{G})
$$
are isomorphisms. On the other hand, the complex
$\check{\mathcal{C}}^\bullet(\mathcal{U}, \mathcal{F})$
is K-flat by Lemma \ref{lemma-K-flat} and we conclude that
$\text{Tot}(
\check{\mathcal{C}}^\bullet(\mathcal{U}, \mathcal{F})
\otimes_A
\check{\mathcal{C}}^\bullet(\mathcal{V}, \mathcal{G}))$
represents the derived tensor product
$R\Gamma(X, \mathcal{F}) \otimes_A^\mathbf{L} R\Gamma(Y, \mathcal{G})$
as claimed.

\medskip\noindent
We still have to show that (\ref{equation-kunneth-on-cech})
is a quasi-isomorphism. We will do this using dimension shifting.
Set $d(\mathcal{F}) = \max \{d \mid H^d(X, \mathcal{F}) \not = 0\}$.
Assume $d(\mathcal{F}) > 0$. Set $U = \coprod\nolimits_{i \in I} U_i$.
This is an affine scheme as $I$ is finite. Denote
$j : U \to X$ the morphism which is the inclusion $U_i \to X$
on each $U_i$. Since the diagonal of $X$ is affine, the morphism
$j$ is affine, see
Morphisms, Lemma \ref{morphisms-lemma-affine-permanence}.
It follows that $\mathcal{F}' = j_*j^*\mathcal{F}$ is $S$-flat, see
Morphisms, Lemma \ref{morphisms-lemma-pushforward-flat-affine}.
It also follows that $d(\mathcal{F}') = 0$ by combining
Cohomology of Schemes, Lemmas
\ref{coherent-lemma-relative-affine-cohomology} and
\ref{coherent-lemma-quasi-coherent-affine-cohomology-zero}.
For all $x \in X$ we have $\mathcal{F}_x  \to \mathcal{F}'_x$
is the inclusion of a direct summand: if $x \in U_i$,
then $\mathcal{F}' \to (U_i \to X)_*\mathcal{F}|_{U_i}$
gives a splitting. We conclude that
$\mathcal{F} \to \mathcal{F}'$ is injective and
$\mathcal{F}'' = \mathcal{F}'/\mathcal{F}$
is $S$-flat as well. The short exact sequence
$0 \to \mathcal{F} \to \mathcal{F}' \to \mathcal{F}'' \to 0$
of flat quasi-coherent $\mathcal{O}_X$-modules
produces a short exact sequence of complexes
$$
0 \to
\text{Tot}(
\check{\mathcal{C}}^\bullet(\mathcal{U}, \mathcal{F})
\otimes_A
\check{\mathcal{C}}^\bullet(\mathcal{V}, \mathcal{G})) \to
\text{Tot}(
\check{\mathcal{C}}^\bullet(\mathcal{U}, \mathcal{F}')
\otimes_A
\check{\mathcal{C}}^\bullet(\mathcal{V}, \mathcal{G})) \to
\text{Tot}(
\check{\mathcal{C}}^\bullet(\mathcal{U}, \mathcal{F}'')
\otimes_A
\check{\mathcal{C}}^\bullet(\mathcal{V}, \mathcal{G})) \to 0
$$
and a short exact sequence of complexes
$$
0 \to
\check{\mathcal{C}}^\bullet(\mathcal{W},
p^*\mathcal{F}
\otimes_{\mathcal{O}_{X \times_S Y}}
q^*\mathcal{G}) \to
\check{\mathcal{C}}^\bullet(\mathcal{W},
p^*\mathcal{F}'
\otimes_{\mathcal{O}_{X \times_S Y}}
q^*\mathcal{G}) \to
\check{\mathcal{C}}^\bullet(\mathcal{W},
p^*\mathcal{F}''
\otimes_{\mathcal{O}_{X \times_S Y}}
q^*\mathcal{G}) \to 0
$$
Moreover, the maps (\ref{equation-kunneth-on-cech}) between these are
compatible with these short exact sequences. Hence it suffices to prove
(\ref{equation-kunneth-on-cech})
is an isomorphism for $\mathcal{F}'$ and $\mathcal{F}''$. Finally,
we have $d(\mathcal{F}'') < d(\mathcal{F})$.
In this way we reduce to the case $d(\mathcal{F}) = 0$.

\medskip\noindent
Arguing in the same fashion for $\mathcal{G}$ we find that we
may assume that both $\mathcal{F}$ and $\mathcal{G}$
have nonzero cohomology only in degree $0$.
Observe that this means that $\Gamma(X, \mathcal{F})$
is quasi-isomorphic to the $K$-flat complex
$\check{\mathcal{C}}^\bullet(\mathcal{U}, \mathcal{F})$
of $A$-modules sitting in degrees $\geq 0$.
It follows that $\Gamma(X, \mathcal{F})$ is a flat $A$-module
(because we can compute higher Tor's against this module
by tensoring with the Cech complex).
Let $V \subset Y$ be an affine open. Consider the affine open covering
$\mathcal{U}_V : X \times_S V = \bigcup_{i \in I} U_i \times_S V$.
It is immediate that
$$
\check{\mathcal{C}}^\bullet(\mathcal{U}, \mathcal{F})
\otimes_A \mathcal{G}(V) =
\check{\mathcal{C}}^\bullet(\mathcal{U}_V,
p^*\mathcal{F} \otimes_{\mathcal{O}_{X \times Y}}
q^*\mathcal{G})
$$
(equality of complexes). By the flatness of $\mathcal{G}(V)$
over $A$ we see that
$\Gamma(X, \mathcal{F}) \otimes_A \mathcal{G}(V) \to
\check{\mathcal{C}}^\bullet(\mathcal{U}, \mathcal{F})
\otimes_A \mathcal{G}(V)$ is a quasi-isomorphism.
Since the sheafification of
$V \mapsto \check{\mathcal{C}}^\bullet(\mathcal{U}_V,
p^*\mathcal{F} \otimes_{\mathcal{O}_{X \times Y}}
q^*\mathcal{G})$ represents
$Rq_*(p^*\mathcal{F} \otimes_{\mathcal{O}_{X \times Y}} q^*\mathcal{G})$
by Cohomology of Schemes, Lemma
\ref{coherent-lemma-separated-case-relative-cech}
we conclude that
$$
Rq_*(p^*\mathcal{F} \otimes_{\mathcal{O}_{X \times Y}} q^*\mathcal{G})
\cong
\Gamma(X, \mathcal{F}) \otimes_A \mathcal{G}
$$
on $Y$ where the notation on the right hand side indicates the module
$$
b^*\widetilde{\Gamma(X, \mathcal{F})} \otimes_{\mathcal{O}_Y} \mathcal{G}
$$
Using the Leray spectral sequence for $q$ we find
$$
H^n(X \times_S Y, p^*\mathcal{F} \otimes_{\mathcal{O}_{X \times Y}}
q^*\mathcal{G}) =
H^n(Y,
b^*\widetilde{\Gamma(X, \mathcal{F})} \otimes_{\mathcal{O}_Y} \mathcal{G})
$$
Using Lemma \ref{lemma-cohomology-base-change} for the morphism
$b : Y \to S = \Spec(A)$ and using that $\Gamma(X, \mathcal{F})$
is $A$-flat we conclude that
$H^n(X \times_S Y, p^*\mathcal{F} \otimes_{\mathcal{O}_{X \times Y}}
q^*\mathcal{G})$ is zero for $n > 0$ and isomorphic to
$H^0(X, \mathcal{F}) \otimes_A H^0(Y, \mathcal{G})$ for $n = 0$.
Of course, here we also use that $\mathcal{G}$ only has
cohomology in degree $0$.
This finishes the proof (except that we should check that the
isomorphism is indeed given by cup product in degree $0$; we omit
the verification).
\end{proof}

\bibliographystyle{plain}
\bibliography{my}
\end{document}
